\documentclass{article}

\usepackage{array}
\usepackage{etoolbox}
\usepackage{fancyhdr}
\usepackage{geometry} 
\usepackage{graphicx}
\usepackage{soul}
\usepackage{titling}

%%%%%%%%%%%%%%%%%%%%%%%%%%%%%%%%%%%%%%%%%%%%%%%%%%%%%%%%%%%%
% BEGIN METADATA: Edit the following as appropriate
%%%%%%%%%%%%%%%%%%%%%%%%%%%%%%%%%%%%%%%%%%%%%%%%%%%%%%%%%%%%

\title{Project Title}  % the title of your project
\newcommand\shorttitle{\thetitle}  % if needed: a shorter title for the document header
% Team members.
\newcommand\firstname{Ifrah Chishti}  % full name
\newcommand\firstid{ic08351}         % ID, e.g. xy01234
\newcommand\secondname{Naaseh Sajid} % full name
\newcommand\secondid{ms08085}        % ID, e.g. xy01234
\newcommand\thirdname{Aina Shakeel}  % full name
\newcommand\thirdid{as08430}         % ID, e.g. xy01234
% Uncomment the rows for the next 2 students if and as needed.
\newcommand\fourthname{Zain Hatim} % full name
\newcommand\fourthid{zh08343}        % ID, e.g. xy01234
% \newcommand\fifthname{Student 5}  % full name
% \newcommand\fifthid{id05}         % ID, e.g. xy01234

%%%%%%%%%%%%%%%%%%%%%%%%%%%%%%%%%%%%%%%%%%%%%%%%%%%%%%%%%%%%
% END METADATA: Do not edit the preamble any further.
%%%%%%%%%%%%%%%%%%%%%%%%%%%%%%%%%%%%%%%%%%%%%%%%%%%%%%%%%%%%

\pagestyle{fancy}
\lhead{Kaavish Proposal}
\chead{\shorttitle}
\rhead{Fall 2025}
\cfoot{Page \thepage}
\renewcommand{\footrulewidth}{0.4pt}

\newenvironment{instruction}{\itshape}{}

\begin{document}

% Cover page.
\begin{titlepage}

\center % Center everything on the page
 
%----------------------------------------------------------------------------------------
%	HEADING SECTIONS
%----------------------------------------------------------------------------------------

\textsc{
  {\LARGE \bf \thetitle}\\\bigskip\bigskip % Your Project Title
  {\large
    Kaavish Project Proposal\\\bigskip
    By}
}\\\bigskip 

%----------------------------------------------------------------------------------------
%	AUTHOR SECTION
%----------------------------------------------------------------------------------------

{\large
  \begin{tabular}{ll}
    \firstname & (\firstid@st.habib.edu.pk) \\
    \secondname & (\secondid@st.habib.edu.pk) \\
    \thirdname & (\thirdid@st.habib.edu.pk) \\
    \ifdef{\fourthname}{\fourthname & (\fourthid@st.habib.edu.pk) \\}{}
    \ifdef{\fifthname}{\fifthname & (\fifthid@st.habib.edu.pk) \\}{}
  \end{tabular}
}
\bigskip\bigskip\bigskip

{\large \today}\\\bigskip\bigskip

\includegraphics[height=5cm]{HU_logo}\\\bigskip
 
%----------------------------------------------------------------------------------------
{\large
  In partial fulfillment of the requirement for \\\medskip
Bachelor of Science \\\medskip
Computer Science
}\\\bigskip\bigskip\bigskip

{\large
  \textsc{
    Dhanani School of Science and Engineering\\\bigskip
    Habib University\\\bigskip 
    Fall 2023
  }\\\bigskip\bigskip 
  Copyright @ 2023 Habib University
}

\end{titlepage}


%%%%%%%%%%%%%%%%%%%%%%%%%%%%%%%%%%%%%%%%%%%%%%%%%%%%%%%%%%%%
% DATA: Populate the rest of the document as instructed.
%%%%%%%%%%%%%%%%%%%%%%%%%%%%%%%%%%%%%%%%%%%%%%%%%%%%%%%%%%%%
\section{Abstract}
\begin{instruction}
    CodeSchool is a global coding literacy platform for children aged 6 - 16, delivering more than 2,000 classes across 20 countries. With a mission to make coding a universal literacy, the organization combines an engaging, game-based curriculum with trained instructors who can inspire the next generation of digital creators. Much like how Uber seamlessly connects riders with drivers, CodeSchool envisions a platform that can intelligently connect students with the right instructor, at the right time, using the right curriculum—whether classes are delivered online or onsite. To achieve this vision at scale, the platform must evolve from a system dependent on manual operations into one that is automated, intelligent, and easy to use.

    Currently, managing the lifecycle of classes at CodeSchool is a complex, time-consuming process. Instructors and administrators rely on spreadsheets to schedule classes, assign students, generate Zoom links, and track progress. Evaluations of teaching quality are also carried out manually, with the team watching recorded sessions and giving scores based on observations. While effective at a smaller scale, this system becomes increasingly difficult to manage as CodeSchool grows, leaving room for scheduling conflicts, inefficiencies, and limited insight into teaching quality. These challenges risk slowing down the organization's ability to expand while maintaining its high standard of student learning outcomes.

    The proposed solution is a web-based platform designed to act as a central instructor portal—streamlining operations, automating repetitive tasks, and enabling data-driven insights. At its core, the platform addresses three key areas: class scheduling, student-instructor matching, and performance evaluation.

    First, the platform will automate class scheduling. By factoring in instructor availability, coding language expertise, and student age and proficiency, it will generate conflict-free schedules. This not only saves administrative effort but also ensures that classes run smoothly without last-minute rescheduling.

    Second, the platform will transform how students are assigned to instructors. Instead of a manual process that relies on human judgment, the system will use a recommendation model inspired by machine learning. Much like matching services in other industries, this model will prioritize placing students with instructors best suited to their learning needs—whether that is based on teaching style, technical strengths, or experience with specific age groups. This will also follow a round robin model to ensure fair distribution of students among instructors, allowing instructors to experience teaching different age groups and fostering a more versatile teaching approach. The result is a more personalized and effective learning experience for every student.

    % Third, to support the logistics of virtual teaching, the platform will integrate with the Zoom API to automatically generate and distribute meeting links for each scheduled class. This eliminates another manual step and ensures that both students and instructors always have seamless access to their sessions.

    Finally, the platform will introduce a more advanced approach to instructor performance analysis. Rather than relying solely on manual reviews, it will leverage automated transcription and analysis of Zoom recordings. Using metrics such as clarity, pacing, student engagement, and adherence to curriculum (just to name a few), the system will provide instructors with actionable feedback. This not only supports professional growth but also ensures that students consistently receive high-quality instruction.

    By bringing these features together into one unified dashboard, the project aims to reduce operational overhead while enabling scalability. Instructors will gain a clear and organized view of their teaching responsibilities, student assignments, and performance metrics. Administrators will benefit from real-time data and analytics that guide decision-making. Most importantly, students will experience smoother, better-matched, and higher-quality instruction—helping CodeSchool continue its mission of making coding a universal literacy.

    In summary, this project is not just a technical upgrade but a step toward reimagining how educational platforms can operate at scale. By blending automation with data-driven insights, it empowers instructors to focus on what they do best—teaching—while the system handles the background logistics. The result is a more efficient, reliable, and impactful platform that supports CodeSchool's global ambitions.
\end{instruction}

\section{Problem definition}
\instruction{
    Currently, managing the lifecycle of classes at CodeSchool is a complex, time-consuming process. Instructors and administrators rely on a spreadsheet based system to schedule classes, assign students, generate Zoom links, and track progress. Evaluations of teaching quality are also carried out manually, with the team watching recorded sessions for each instructor and giving scores based on observations. While effective at a smaller scale, this system becomes increasingly difficult to manage as CodeSchool grows, leaving room for scheduling conflicts, inefficiencies, and limited insight into teaching quality. These challenges risk slowing down the organization's ability to expand while maintaining its high standard of student learning outcomes.
}

\section{Social relevance}
\instruction{
    The project addresses the broader societal challenge of ensuring equitable access to quality education in an increasingly digital world. As demand for coding literacy grows globally, many children—particularly in rural areas—lack access to qualified instructors and structured learning opportunities. Current manual systems of scheduling, matching, and evaluation are not scalable, limiting the reach and effectiveness of coding education platforms. By automating these processes and introducing data-driven insights, the project not only streamlines operations but also ensures that students are consistently matched with the right instructors and receive high-quality teaching. In doing so, it supports the mission of making coding a universal literacy, reduces barriers to access, and models how education systems can expand sustainably to meet the needs of diverse learners worldwide.
}

\section{Originality/Novelty}
\instruction{
    Solving this problem creates significant value for both learners and educators. By automating class scheduling, student–instructor matching, and performance evaluation, the platform removes operational bottlenecks that traditionally slow down digital education programs. Instructors gain time and clarity to focus on teaching rather than administrative work, while students benefit from being paired with instructors who best fit their learning needs—improving engagement, consistency, and outcomes. For administrators, the centralized dashboard provides real-time visibility into scheduling, payments, and teaching quality, supporting more informed decision-making and sustainable growth.

    Compared with existing solutions, which often rely on spreadsheets, messaging apps, or generic learning management systems, this platform is purpose-built for coding education at scale. Tools like Zoom and Google Classroom are excellent for video delivery and content sharing, but they do not address the deeper challenge of intelligently matching students to instructors or evaluating teaching quality. Likewise, traditional CRMs or scheduling tools may handle logistics but lack an educational focus and offer no integration with performance analytics. The added value of this project lies in combining all these elements—matching, scheduling, conferencing, and evaluation—into a single, seamless experience. Current solutions operate in silos, but this platform brings them together into one cohesive ecosystem, making it uniquely positioned to scale and sustain quality learning worldwide.
}

\section{CS Contribution}
The project integrates concepts from multiple areas of Computer Science, directly reflecting the knowledge and skills gained through the courses offered at Habib. Key contributions include:

\begin{itemize}
    \item \textbf{Deep Learning:} Applying advanced machine learning techniques to evaluate instructor performance. Models are trained on audio and text data using frameworks such as TensorFlow and PyTorch, leveraging Natural Language Processing (NLP), sentiment analysis, and speech recognition to generate meaningful effectiveness scores.
    
    \item \textbf{Algorithm Design and Analysis:} Developing an intelligent scheduling engine that uses optimization techniques, such as greedy algorithms and genetic algorithms, to minimize conflicts and maximize student--instructor matching efficiency. Algorithmic complexity analysis ensures the system remains efficient at scale.
    
    \item \textbf{Data Structures and Algorithms:} Implementing efficient data management through appropriate structures (heaps, graphs, hash maps) to handle scheduling, student allocation, and query processing. These choices ensure low-latency responses and robust handling of large datasets.
    
    \item \textbf{Web and Mobile Development:} Designing and building a scalable, user-friendly application accessible across devices. This includes developing modular front-end and back-end architectures, integrating APIs (Zoom, scheduling, billing), and ensuring usability through responsive design and intuitive dashboards for instructors, students, and administrators.
    
    \item \textbf{Software Engineering:} Developing a modular, scalable web application that integrates multiple APIs (Zoom, scheduling, billing). Emphasis is placed on maintainability, usability, and meeting non-functional requirements such as reliability and availability.
\end{itemize}

Collectively, these contributions demonstrate the application of theoretical knowledge to solve real-world problems, bridging advanced algorithms, deep learning, and scalable software engineering practices.



\section{Scope and Deliverables}
Given the year-long duration and a team of four members, the project scope is quite feasible considering how we will all be working according to our strengths. The system is divided into manageable modules that allow development to continue side by side instead of waiting for each module to finish. The scope includes:

\begin{itemize}
    \item Development of an instructor portal with availability management and performance analytics.
    \item Implementation of a scheduling algorithm that assigns classes automatically while preventing conflicts.
    \item Creation of a student--instructor matching algorithm leveraging computational intelligence.
    \item Integration with the Zoom API for automated link generation and distribution.
    \item Machine learning models for instructor performance analysis based on audio and text data from class recordings.
    \item Administrative dashboard providing reporting, analytics, and payment processing.
\end{itemize}

\textbf{Deliverables:}
\begin{itemize}
    \item Functional web-based platform with separate dashboards for instructors and administrators.
    \item Database schema and implementation for managing schedules, profiles, and performance records.
    \item Smart scheduling engine with optimization-based algorithms (this will include the relevant ML models as well).
    \item Documentation including design artifacts and API specifications.
    \item Deployment-ready system demonstration.
\end{itemize}


\section{Feasibility}
\instruction{List the resources, e.g. datasets, compute resources, software libraries, hardware, required for the project. Mention how you expect to access and utilize them for the project.}

\section{Team dynamics}
The project team consists of four members, each bringing complementary expertise across backend development, frontend design, databases, and machine learning. This distribution of roles ensures that the interdisciplinary requirements of the system are well-covered and integration can be achieved smoothly.

\begin{itemize}
    \item \textbf{Naaseh Sajid (ms08085):} Specializes in backend development, API integration, and database management. Skilled in Node.js, Next.js, Ruby on Rails, React, Python, and PostgreSQL. He is responsible for developing the backend services that connect the web application with the database and evaluation modules, as well as maintaining and populating the database.
    
    \item \textbf{Ifrah Chishti (ic08351):} Contributes to both frontend and backend development, with additional expertise in databases and machine learning. Experienced in Node.js, Python, PostgreSQL, and data science. She is primarily responsible for implementing design and code logic aligned with CodeSchool’s requirements, while also providing backend support alongside Naaseh and Aina.
    
    \item \textbf{Aina Shakeel (as08430):} Focuses on machine learning and computational intelligence. Proficient in Python, TensorFlow, Scikit-learn, and deep learning techniques. She is tasked with developing and integrating the student–instructor matching model, and also supports backend architecture to enable robust data collection for future performance analytics.
    
    \item \textbf{Zain Hatim (zh08343):} Specializes in frontend design and user experience. Skilled in React, React Native, Node.js, and Python. He is responsible for building the frontend interface in React, ensuring usability, simplicity, and an aesthetically appealing design, as well as addressing any UI/UX issues that arise.
\end{itemize}


\section{Tech stack}
The project employs a modern and scalable technology stack, chosen for its robustness and ability to support modular development.

\begin{itemize}
    \item \textbf{Frontend:} React.js / Next.js for building responsive, component-based UIs with smooth state management and routing.
    \item \textbf{Backend:} Ruby on Rails or Next.js (API routes) to implement RESTful services for scheduling, user management, and system logic.
    \item \textbf{Machine Learning:} Python with libraries such as TensorFlow, PyTorch, and SpeechRecognition for audio feature extraction, NLP-based analysis, and building instructor performance models.
    \item \textbf{Database:} PostgreSQL as the primary relational database for storing instructor availability, student profiles, class schedules, and performance metrics.
    \item \textbf{Video Conferencing Integration:} Zoom API for automatic meeting creation and link distribution; Google Meet as a fallback option.
    \item \textbf{Deployment and Version Control:} Git/GitHub for version control, Docker for containerization, and cloud-based hosting (e.g., AWS or Heroku) for deployment and scalability.
\end{itemize}


\section{References}
\instruction{List your references.}

% External advisor undertaking.
\begin{titlepage}

  \newcolumntype{C}[1]{>{\centering\arraybackslash\hspace{0pt}}p{#1}}

  \centerline{\textbf{\ul{Undertaking of Kaavish advisement as an External Supervisor}}}
  \bigskip\bigskip

  I hereby affirm that I have read the project details as described on the preceding pages and agree to undertake advisement of this Kaavish project as an External Supervisor. I understand that this role entails the following.
  \begin{description}
  \item[Meeting] Meeting the project team regularly, at least once every two weeks, for the entire duration of the Kaavish. The meetings may be held remotely if required.
  \item[Advisement] Providing supervision and advice to the team in order to ensure steady progress of the project toward its goals.
  \item[Liaison] Liaising with the Internal Supervisor as required, e.g. to provide feedback or engage in grading.
  \item[Other] Any other task, depending on availability and suitability, relevant to the Kaavish as communicated by the Internal Supervisor or Kaavish Working Group.
  \end{description}

  \bigskip\bigskip\bigskip
  
  \noindent
  \begin{tabular}{@{}lp{.4\textwidth}}
    Name: & \hrulefill\\\\
    Email: & \hrulefill\\\\
    Phone: & \hrulefill\\\\
    Designation: & \hrulefill\\\\
    Affiliation: & \hrulefill\\\\\\
    Signature: & \hrulefill\\
  \end{tabular}
\end{titlepage}


\end{document}

%%% Local Variables:
%%% mode: latex
%%% TeX-master: t
%%% End:
